%
%  chapter4.tex — Implementação
%
\chapter{Implementação}
\label{cha:implementacao}

\section{Decisões Arquiteturais}
\label{sec:decisoes_arquiteturais}

O processo de desenvolvimento do \textit{Play4Change} foi suportado por um conjunto de Registos de Decisão Arquitetural (\textit{Architecture Decision Records} --- ADRs), que documentam formalmente as escolhas técnicas mais significativas, os critérios de avaliação considerados e as alternativas rejeitadas. Este processo garante rastreabilidade e facilita a evolução fundamentada do sistema.

As principais decisões estão detalhadas no Apêndice~\ref{cha:adrs}. De forma sumária, destacam-se:

\begin{itemize}
  \item \textbf{Gestão de erros}: adoção do padrão \textit{Arrow Either} com DSL \textit{Raise}, que permite modelar explicitamente os ramos de erro no domínio da aplicação sem recorrer a exceções;
  \item \textbf{Armazenamento vetorial}: extensão \textit{pgvector} sobre PostgreSQL, evitando a introdução de um serviço adicional dedicado (e.g., Qdrant, Weaviate);
  \item \textbf{Mensageria}: eventos de domínio Spring, internos aos contextos delimitados, sem necessidade de um \textit{broker} externo (e.g., Kafka, RabbitMQ);
  \item \textbf{Armazenamento de ficheiros}: MinIO (serviço Docker compatível com S3), com migração para AWS S3 prevista em fases ulteriores;
  \item \textbf{Implantação}: Docker Compose sobre VPS, com critérios explícitos de migração para Kubernetes definidos no ADR-009;
  \item \textbf{Estado no cliente móvel}: padrão MVI (\textit{Model-View-Intent}) via \textit{BaseComponent} e \textit{Decompose MutableValue}.
\end{itemize}

\section{Stack Tecnológico}
\label{sec:stack}

A Tabela~\ref{tab:stack} apresenta o \textit{stack} tecnológico adotado em cada camada da plataforma.

\begin{table}[ht]
  \caption{Stack tecnológico do Play4Change}
  \label{tab:stack}
  \vskip0.5em
  \centering
  \begin{tabular}{ll}
    \toprule
    \textbf{Camada} & \textbf{Tecnologias} \\
    \midrule
    Cliente móvel       & Kotlin Multiplatform, Compose Multiplatform, Decompose \\
    Cliente web         & \textit{(a definir)} \\
    Servidor            & Kotlin, Spring Boot, REST, PostgreSQL, pgvector \\
    Agente de IA        & LangChain4j, Mistral (LLM), pgvector (embeddings) \\
    Autenticação        & Magic Link (SHA-256, Resend API), Google OAuth, Facebook OAuth \\
    Armazenamento       & PostgreSQL, MinIO (S3-compatible) \\
    Observabilidade     & Micrometer, Prometheus, Grafana \\
    CI/CD               & GitHub Actions \\
    Implantação         & Docker Compose, VPS \\
    \bottomrule
  \end{tabular}
\end{table}

\section{Estrutura Modular do Servidor}
\label{sec:modulos_servidor}

O servidor foi organizado segundo os princípios de arquitetura por contextos delimitados (\textit{Bounded Contexts}), em alinhamento com os princípios de \textit{Domain-Driven Design} (DDD)~\cite{evans2003ddd}. Os módulos principais são:

\begin{description}
  \item[\texttt{:domain}] Entidades de domínio, interfaces de repositório e lógica de negócio central;
  \item[\texttt{:application}] Casos de uso e serviços de orquestração;
  \item[\texttt{:infrastructure}] Adaptadores de persistência (PostgreSQL, MinIO), clientes externos e configuração;
  \item[\texttt{:api}] Controladores REST, DTOs e mapeamentos;
  \item[\texttt{:ai-agent:langchain}] Módulo Gradle dedicado à integração LangChain4j/Mistral e à lógica de adaptação.
\end{description}

Esta separação assegura que a lógica de domínio permanece independente dos detalhes de infraestrutura, facilitando a testabilidade e a evolução do sistema. As dependências entre módulos estão ilustradas na Figura~\ref{fig:modulos}.

\begin{figure}[htbp]
  \centering
  % TODO: substituir por img/modulos.pdf quando disponível
  \begin{tikzpicture}[
    >=Stealth, font=\small,
    mod/.style={draw, rounded corners=4pt, fill=blue!12,
                minimum width=3.8cm, minimum height=0.8cm,
                align=center, font=\small\ttfamily},
    aimod/.style={mod, fill=orange!18}
  ]
  % Module nodes
  \node[mod]   (dom) at (5.0,  0)   {:domain};
  \node[mod]   (app) at (5.0, -2.2) {:application};
  \node[mod]   (inf) at (1.5, -4.4) {:infrastructure};
  \node[mod]   (api) at (5.0, -4.4) {:api};
  \node[aimod] (ai)  at (8.5, -4.4) {:ai-agent:langchain};
  % Dependency arrows (arrow = depends on)
  \draw[->] (app) -- (dom)  node[midway,right,font=\tiny] {depende de};
  \draw[->] (inf) -- (app);
  \draw[->] (api) -- (app);
  \draw[->] (ai)  -- (app);
  \draw[->] (inf) to[out=90,in=210] (dom);
  \draw[->] (api) to[out=90,in=330] (dom);
  \draw[->] (ai)  to[out=90,in=315] (dom);
  \end{tikzpicture}
  \caption{Estrutura modular do servidor e dependências entre módulos Gradle}
  \label{fig:modulos}
\end{figure}

\section{Desenho da API REST}
\label{sec:api_rest}

A API REST do servidor está organizada por recurso, seguindo as convenções de nomenclatura e os métodos HTTP da arquitetura REST~\cite{fielding2000rest}. A Tabela~\ref{tab:api_endpoints} apresenta os principais grupos de \textit{endpoints}.

\begin{table}[ht]
  \caption{Principais grupos de \textit{endpoints} da API REST}
  \label{tab:api_endpoints}
  \vskip0.5em
  \centering
  \begin{tabular}{lp{2.8cm}p{6.5cm}}
    \toprule
    \textbf{Prefixo} & \textbf{Métodos} & \textbf{Função} \\
    \midrule
    \texttt{/auth}          & POST             & Autenticação: \textit{magic link}, Google OAuth, Facebook OAuth \\
    \texttt{/users}         & GET, PATCH       & Gestão do perfil do utilizador autenticado \\
    \texttt{/topics}        & GET, POST, PATCH, DELETE & Gestão de tópicos (requer papel de administrador) \\
    \texttt{/tasks}         & GET, POST        & Obtenção de tarefas diárias e submissão de respostas \\
    \texttt{/badges}        & GET              & Consulta de \textit{badges} disponíveis e obtidos \\
    \texttt{/progress}      & GET              & Histórico de progresso do cidadão \\
    \texttt{/adaptive}      & GET              & Ramo adaptativo ativo do cidadão \\
    \texttt{/admin/metrics} & GET              & Métricas de utilização (requer papel de administrador) \\
    \bottomrule
  \end{tabular}
\end{table}

A especificação completa da API encontra-se no contrato OpenAPI gerado automaticamente em tempo de execução, disponível em \texttt{/v3/api-docs}.

\section{Agente de IA: Implementação}
\label{sec:ia_implementacao}

O agente de inteligência artificial é implementado no módulo \texttt{:ai-agent:langchain}, utilizando a \textit{framework} LangChain4j para orquestrar as interações com o modelo de linguagem Mistral. O fluxo de geração de conteúdo segue um \textit{pipeline} estruturado:

\begin{enumerate}
  \item O administrador define um tópico e os seus objetivos de aprendizagem no portal de administração;
  \item O agente recebe o pedido de geração e constrói um \textit{prompt} estruturado com os parâmetros do tópico;
  \item O modelo Mistral gera as tarefas de aprendizagem (questões, explicações, exercícios);
  \item O conteúdo gerado é validado, persistido e associado ao tópico correspondente;
  \item Os embeddings semânticos do conteúdo são calculados e armazenados em \textit{pgvector} para recuperação futura.
\end{enumerate}

O mecanismo de deteção de dificuldades monitoriza, em tempo real, métricas de desempenho do cidadão (e.g., taxa de erros, tempo de resposta) e, quando o limiar de dificuldade é excedido, efetua uma pesquisa semântica no repositório \textit{pgvector} para identificar percursos adaptativos previamente bem-sucedidos com perfis de dificuldade semelhantes. O \textit{pipeline} completo de geração de conteúdo encontra-se ilustrado na Figura~\ref{fig:pipeline_ia}.

\begin{figure}[htbp]
  \centering
  % TODO: substituir por img/pipeline_ia.pdf quando disponível
  \begin{tikzpicture}[
    >=Stealth, font=\small,
    pbox/.style={draw, rounded corners=3pt, fill=blue!12,
                 minimum width=4.0cm, minimum height=0.8cm, align=center},
    store/.style={pbox, fill=orange!18},
    ext/.style={pbox, dashed, fill=green!12}
  ]
  % Main pipeline (vertical)
  \node[pbox]  (admin)  at (0, 0)    {Administrador define tópico};
  \node[pbox]  (prompt) at (0,-1.8)  {Construção do prompt\\(objetivos pedagógicos)};
  \node[ext]   (llm)    at (0,-3.6)  {Mistral API\\\scriptsize (geração de conteúdo)};
  \node[pbox]  (valid)  at (0,-5.4)  {Validação do conteúdo gerado};
  \node[store] (pers)   at (0,-7.2)  {Persistência (PostgreSQL)};
  \node[store] (emb)    at (0,-9.0)  {Cálculo de embeddings\\\scriptsize (LangChain4j)};
  \node[store] (pgv)    at (0,-10.8) {Armazenamento em pgvector};
  % Arrows
  \foreach \a/\b in {admin/prompt, prompt/llm, llm/valid, valid/pers, pers/emb, emb/pgv}
    \draw[->] (\a) -- (\b);
  % Side branch: reuse
  \node[pbox, fill=gray!15, text width=3.8cm, align=center] (reuse)
    at (6.5,-10.8) {Reutilização:\\pesquisa por semelhança\\(kNN em perfis de dificuldade)};
  \draw[<->, dashed] (pgv.east) -- (reuse.west)
    node[midway,above,font=\tiny] {kNN search};
  \end{tikzpicture}
  \caption{Pipeline de geração de conteúdo e reutilização do agente de IA}
  \label{fig:pipeline_ia}
\end{figure}

\section{Observabilidade}
\label{sec:observabilidade}

A plataforma integra um sistema de observabilidade baseado em Micrometer, Prometheus e Grafana, com sete métricas de domínio personalizadas que permitem monitorizar indicadores chave do negócio, como: taxa de conclusão de microcompetências, frequência de ativação de percursos adaptativos, distribuição de \textit{badges} atribuídos e latência do agente de IA. Os \textit{dashboards} Grafana são provisionados a partir do código-fonte, garantindo reprodutibilidade e rastreabilidade da configuração de monitorização.

\section{Segurança}
\label{sec:seguranca}

A segurança da plataforma foi auditada formalmente (ADR-016), tendo sido identificadas e corrigidas onze vulnerabilidades. A única lacuna por resolver — denominada G4 — refere-se à validação de audiência em fluxos Facebook OAuth e encontra-se documentada e calendarizada para resolução em iterações subsequentes.

\section{Diagrama de Implantação}
\label{sec:diagrama_implantacao}

A plataforma é implantada através de Docker Compose numa VPS, com os seguintes serviços orquestrados:

\begin{description}
  \item[\texttt{app}] Servidor Spring Boot (porta 8080); configurado por variáveis de ambiente para base de dados, serviços externos e chaves de autenticação;
  \item[\texttt{db}] PostgreSQL~16 com extensão \textit{pgvector}, acessível internamente na porta 5432;
  \item[\texttt{minio}] Armazenamento compatível com S3 (portas 9000/9001), utilizado para ficheiros estáticos e imagens de \textit{badges};
  \item[\texttt{prometheus}] Recolha de métricas via \textit{endpoint} Micrometer exposto pelo servidor (porta 9090);
  \item[\texttt{grafana}] Visualização de \textit{dashboards} (porta 3000), provisionados a partir do código-fonte.
\end{description}

Os serviços externos --- API Mistral, API Resend, Google OAuth e Facebook OAuth --- são consumidos pelo servidor via HTTPS, sem exposição na rede interna do Docker Compose.

\begin{figure}[htbp]
  \centering
  % TODO: substituir por img/implantacao.pdf quando disponível
  \resizebox{\textwidth}{!}{%
  \begin{tikzpicture}[
    >=Stealth, font=\small,
    svc/.style={draw, rounded corners=3pt, fill=blue!12,
                minimum width=3.0cm, minimum height=1.1cm,
                align=center, text width=2.8cm},
    ext/.style={svc, fill=green!12, dashed},
    usr/.style={svc, fill=gray!15}
  ]
  % VPS outer box
  \draw[thick, rounded corners=8pt, fill=white] (-0.8,-8.5) rectangle (14.5,2.0);
  \node[font=\small\bfseries] at (6.8, 1.7) {VPS};
  % Docker Compose inner box
  \draw[dashed, rounded corners=5pt, fill=blue!4] (0,-8.0) rectangle (12.5,1.2);
  \node[font=\footnotesize\itshape, gray] at (6.25, 0.95)
    {Rede Docker Compose (interna)};
  % Services top row
  \node[svc] (app)  at (2.0, 0.0) {\textbf{app}\\\scriptsize Spring Boot\\:8080};
  \node[svc] (db)   at (6.2, 0.0) {\textbf{db}\\\scriptsize PostgreSQL 16\\+ pgvector~:5432};
  \node[svc] (minio)at (10.5,0.0) {\textbf{minio}\\\scriptsize MinIO S3\\:9000/:9001};
  % Services bottom row
  \node[svc] (prom) at (4.0,-3.5) {\textbf{prometheus}\\\scriptsize Métricas\\:9090};
  \node[svc] (grf)  at (8.5,-3.5) {\textbf{grafana}\\\scriptsize Dashboards\\:3000};
  % Internal connections
  \draw[<->] (app.east)  -- (db.west);
  \draw[<->] (app.east)  -- (minio.west);
  \draw[->]  (app.south) to[bend right=10] (prom.north);
  \draw[->]  (prom.east) -- (grf.west);
  % External services (below VPS)
  \node[ext] (mis) at (1.5, -11.0) {\textbf{Mistral API}\\\scriptsize LLM};
  \node[ext] (res) at (5.0, -11.0) {\textbf{Resend}\\\scriptsize Email};
  \node[ext] (gg)  at (8.5, -11.0) {\textbf{Google OAuth}\\\scriptsize JWKS};
  \node[ext] (fb)  at (12.0,-11.0) {\textbf{Facebook OAuth}\\\scriptsize Graph API};
  % App to external (HTTPS)
  \foreach \e in {mis,res,gg,fb}
    \draw[->,dashed] (app.south) to[out=270,in=90] (\e.north);
  \node[font=\tiny,gray] at (6.8,-9.5) {HTTPS};
  % Users (right side)
  \node[usr] (cit) at (17.5, 0.5) {\textbf{Cidadão}\\\scriptsize App Android/iOS};
  \node[usr] (adm) at (17.5,-2.5) {\textbf{Administrador}\\\scriptsize Portal Web};
  \draw[->] (cit.west) -- (app.east) node[midway,above,font=\tiny] {HTTPS};
  \draw[->] (adm.west) to[out=180,in=0] (app.east);
  \end{tikzpicture}}
  \caption{Diagrama de implantação do Play4Change (Docker Compose sobre VPS)}
  \label{fig:implantacao}
\end{figure}
